\documentclass{article}
\usepackage[UTF8]{ctex}
\usepackage{amsmath,amssymb,amsfonts}
\usepackage{enumitem}
\usepackage{geometry}
\geometry{a4paper, margin=1in}
% === 错题本排版优化（不参与题目标识） ===
\setlength{\parindent}{0pt}
\setlength{\parskip}{0.35em plus 0.10em minus 0.08em}
\setlength{\emergencystretch}{3em}
\setlist{itemsep=0.25em, topsep=0.40em, parsep=0pt}
\allowdisplaybreaks[2]
% === 错题本排版优化结束 ===
\begin{document}
\section*{武汉大学2025年高等代数真题}
\section*{122．武汉大学2025年高等代数试卷}

\subsection{题目1}
1．已知复矩阵 \(A=\left(\begin{array}{ccc}0 & 0 & 0 \\ 13 & 0 & 0 \\ 1 & 15 & 0\end{array}\right)\) ，证明：不存在三阶复矩阵 \(X\) ，使得矩阵方程 \(X^{2}=A\) 有解．

\subsection{题目2}
2．\(A=\left(a_{i j}\right)_{3 \times 3}\) ，\(f\) 是线性映射，满足

\[
\begin{gathered}
f: M_{3 \times 3}(\mathbb{R}) \rightarrow \mathbb{R} \\
\left(a_{i j}\right)_{3 \times 3} \mapsto 7 a_{23}
\end{gathered}
\]

是否存在 \(3 \times 3\) 矩阵 \(B\) ，使得 \(f(A)=\operatorname{tr}\left(A B^{\mathrm{T}}\right)\) 是任意的 \(A \in M_{3 \times 3}(\mathbb{R})\) 成立？其中 \(\operatorname{tr}\) 表示矩阵的迹．

\subsection{题目3}
3．已知 \(A=\left(\begin{array}{ccc}4 & 8 & -5 \\ 0 & 2 & -1 \\ 2 & 1 & -1\end{array}\right)\) ，线性变换

\[
\begin{aligned}
\sigma: M_{3 \times 3}(\mathbb{R}) & \rightarrow M_{3 \times 3}(\mathbb{R}) \\
B & \mapsto A B
\end{aligned}
\]

（1）求 \(\sigma\) 核空间的一组基。\\
（2）求 \(\sigma\) 像空间的一组基。

\subsection{题目4}
4．矩阵 \(A=\left(\begin{array}{rrrrrrr}2 & 1 & & & & & \\ 1 & 4 & 1 & & & & \\ & 1 & 6 & & & & \\ & & \ddots & \ddots & 1 & & \\ & & & 1 & 2(n-1) & 1 \\ & & & & 1 & & 2 n\end{array}\right)\)\\
（1）若矩阵 \(B\) 满足 \(A B=B A\) ，证明：\(B\) 可对角化．\\
（2）若 \(C(A)=\{B \mid A B=B A\}\) ，求 \(C(A)\) 的维数．

\subsection{题目5}
5．已知矩阵 \(A=\left(\begin{array}{ccc}3 & \frac{a}{2} & -\frac{3}{2} \\ 0 & a & 21 \\ 1 & \frac{3}{14} & 5\end{array}\right)\) 的特征多项式有重根，问 \(A\) 是否可对角化？

\subsection{题目6}
6．已知 \(A=\left(\begin{array}{rrrr}0 & & & -1 \\ 2 & 0 & & -2 \\ 2 & 0 & -3 \\ & \ddots & \ddots & \vdots \\ & & 2 & -n\end{array}\right)\) ，问 \(A\) 是否可逆？若 \(A\) 可逆，求出 \(A^{-1}\) ．

\subsection{题目7}
7．\(V\) 是数域 \(\mathbb{R}\) 上的 6 维线性空间， \(\mathcal{A}\) 是 \(V\) 上的线性变换，实系数多项式 \(f(x)=\) \(\left(x^{2}+2\right)\left(x^{2}+3\right)\left(x^{2}+5\right)\) 为 \(\mathcal{A}\) 的特征多项式，证明：\(V\) 可表示为三个 \(\mathcal{A}-\) 不变子空间的直和。

\subsection{题目8}
8．\(A, B\) 都是 \(3 \times 3\) 的复矩阵。

\begin{itemize}
  \item[（1）] 若 \(A B-B A=A\) ，证明：\(A\) 是幂零矩阵．
  \item[（2）] 线性变换
\[
\begin{gathered}
\tau: M_{3 \times 3}(\mathbb{C}) \rightarrow M_{3 \times 3}(\mathbb{C}) \\
C \mapsto A C-C A
\end{gathered}
\]
其中 \(A\) 满足 \(A B-B A=A\) ，证明：\(\tau\) 是幂零变换．
\end{itemize}
\end{document}
